% Источники: exam/materials/моделирование.pdf, раздел 1; GitHub HumsterProgrammer/iu9-notes/8sem/Моделирование/Моделирование_лаб-3_06-04-26.md.
\section{Построение логической модели. СДНФ и СКНФ.}

\subsection*{Основные определения}

Логическая модель задаётся \textbf{таблицей истинности}: для каждого набора значений бинарных входных признаков $x_1,\ldots,x_k \in \{0,1\}$ указывается значение выходного признака $y \in \{0,1\}$.

\begin{itemize}
  \item \textbf{СДНФ} (совершенная дизъюнктивная нормальная форма) — дизъюнкция (OR) всех \emph{минтермов}: конъюнкций (AND), обращающихся в 1 на строках с $y=1$.
  \item \textbf{СКНФ} (совершенная конъюнктивная нормальная форма) — конъюнкция (AND) всех \emph{макстермов}: дизъюнкций (OR), обращающихся в 0 на строках с $y=0$.
\end{itemize}

\subsection*{Алгоритм построения}

\begin{enumerate}
  \item \textbf{Бинаризовать} входные признаки (например, по медиане выборки).
  \item Составить \textbf{таблицу истинности}: перебрать все $2^k$ комбинаций, для каждой определить $y$ (по рабочей части выборки).
  \item \textbf{СДНФ}: для каждой строки с $y=1$ выписать конъюнкцию вида
    $x_1^{(a_1)} \land x_2^{(a_2)} \land \cdots \land x_k^{(a_k)}$,
    где $x_j^{(1)} = x_j$, $x_j^{(0)} = \overline{x_j}$.
    СДНФ — дизъюнкция всех таких конъюнкций.
  \item \textbf{СКНФ}: для каждой строки с $y=0$ выписать дизъюнкцию вида
    $x_1^{(b_1)} \lor x_2^{(b_2)} \lor \cdots \lor x_k^{(b_k)}$,
    где $x_j^{(0)} = x_j$, $x_j^{(1)} = \overline{x_j}$ (знаки \emph{инвертированы} по сравнению с СДНФ).
    СКНФ — конъюнкция всех таких дизъюнкций.
  \item \textbf{Верификация}: проверить совпадение СДНФ и СКНФ на контрольной части выборки.
\end{enumerate}

\subsection*{Пример из лабораторного конспекта}

Пусть есть три признака $x_1,x_2,x_3 \to y$. После бинаризации перебираются все возможные варианты:

\begin{center}
\begin{tabular}{ccc|c}
  \hline
  $x_1$ & $x_2$ & $x_3$ & $y$ \\
  \hline
  0 & 0 & 0 & 1 \\
  0 & 0 & 1 & 1 \\
  0 & 1 & 0 & 0 \\
  0 & 1 & 1 & 1 \\
  1 & 0 & 0 & 0 \\
  1 & 0 & 1 & 1 \\
  1 & 1 & 0 & 0 \\
  1 & 1 & 1 & 1 \\
  \hline
\end{tabular}
\end{center}

\textbf{СДНФ} (строки с $y=1$):
\[
  \text{СДНФ} =
  \overline{x_1}\overline{x_2}\overline{x_3}
  \lor \overline{x_1}\overline{x_2}x_3
  \lor \overline{x_1}x_2x_3
  \lor x_1\overline{x_2}x_3
  \lor x_1x_2x_3.
\]

\textbf{СКНФ} (строки с $y=0$):
\[
  \text{СКНФ} =
  (x_1 \lor \overline{x_2} \lor x_3)
  \land (\overline{x_1} \lor x_2 \lor x_3)
  \land (\overline{x_1} \lor \overline{x_2} \lor x_3).
\]

Логическую модель строят по рабочей выборке; затем по рабочей выборке проверяют модель. Выборку разбивают на рабочую и контрольную, упорядочивая лексикографически по бинаризации.

\clearpage
